\RequirePackage{silence}
\WarningsOff % Turn off all warnings from all packages
\ErrorsOff   % (Caution) Use only if you are sure the code has no severe logic errors

\documentclass[a4paper,14pt, twoside]{memoir}
\usepackage[utf8x]{inputenc}
\usepackage[english]{babel}    
\usepackage{url}
\usepackage{microtype}
\usepackage{color}
\usepackage{amsmath}
\usepackage{amsthm}         % Provides proof environment
\usepackage{amssymb}
\usepackage{mathtools}     % Includes amsmath but optimized and with more features
\usepackage{hyperref}
\usepackage{enumitem}
\usepackage{autobreak}
\usepackage{array}          % Support for notation table formatting
\usepackage{booktabs}       % Professional table rules

\hypersetup{
	colorlinks=true,
	linkcolor=blue,
	filecolor=magenta,
	urlcolor=cyan
}

% Theorem Environments (Đánh số theo Chapter)
\newtheorem{definition}{Definition}[chapter]
\newtheorem{md}{Proposition}[chapter]
\newtheorem{theorem}{Theorem}[chapter]
\newtheorem{proposition}{Proposition}[chapter]
\newtheorem{corollary}{Corollary}[chapter]
\newtheorem{remark}{Remark}[chapter]
\newtheorem{example}{Example}[chapter]

% Declare Math Operators
\DeclareMathOperator{\Tr}{Tr}
\DeclareMathOperator{\Sym}{Sym}
\DeclareMathOperator{\setu}{s}
\DeclareMathOperator{\dimU}{\dim}
\DeclareMathOperator{\w}{w}
\DeclareMathOperator{\card}{card}
\DeclareMathOperator{\pr}{pr}
\DeclareMathOperator{\pgr}{pgr}
\DeclareMathOperator{\scale}{scale}
\DeclareMathOperator{\rad}{rad}
\DeclareMathOperator{\Val}{Val}
\DeclareMathOperator{\Conf}{Conf}

% Custom commands
\newcommand{\set}[1]{\left\{#1\right\}} % set
\newcommand{\tbd}[1]{\mathcal{#1}} % uncertain set
\newcommand{\s}[1]{\set{#1}_\mathrm{u}} % uncertain number
\newcommand{\kbd}[1]{{#1}_\mathrm{u}} % uncertain interval
\newcommand{\bd}[1]{\left(#1\right)_\mathrm{u}} % uncertain expression
\newcommand{\tsbd}[1]{\mathbf{U}(\mathbb{#1})}
\newcommand{\tsbdcapn}[2]{\mathbf{U^{#2}}(\mathbb{#1})}
\newcommand{\tsbdts}[1]{\mathbf{U}_{\w}(\mathbb{#1})} % set of uncertain numbers
\newcommand{\ddclmd}{\mathcal{V}} % truth measure of a proposition
\newcommand{\ts}[1]{\w_{#1}} % Weight
\newcommand{\gtcl}[2]{\left(#1\right)_{#2}} % truth measure
\newcommand{\gr}[1]{\text{gr}(#1)} % graph of a function
\newcommand{\hpw}[2]{#1_1\left(#2\right)} % point-wise function
\newcommand{\hbd}[2]{#1_m\left(#2\right)} % uncertain function
\newcommand{\eq}[1]{\left[#1\right]_{eq}} % eqs operator
\newcommand{\kgdd}[1]{\left(#1\right)_1} % pointwise operator
\newcommand{\kgddmr}[1]{\left(#1\right)_{1'}} % real operator
\newcommand{\kgdt}[1]{\left(#1\right)_m} % pointwise operator
\newcommand{\kgdtmr}[1]{\left(#1\right)_{m'}} % pointwise operator

\title{\textbf{Scenario-Based Uncertain Number Theory}}
\author{
	Tran Manh Cuong\thanks{Email: \href{mailto:tmcuong2408@gmail.com}{tmcuong2408@gmail.com} | Phone: 0353237140} \\[0.5em]
	\small Alumnus of Faculty of Mathematics and Informatics (K11)
	\\ Thai Nguyen University of Sciences \\
	\small Address: DJ7 Street, Thoi Hoa, Ho Chi Minh City, Vietnam \\
	\small \textit{Research Field: Mathematics of Uncertainty}
}
\date{Last updated: \today}

\begin{document}
	
	\maketitle
	
	\begin{abstract}
		This work reformulates classical Number Theory into the framework of Uncertain Mathematics using weak binary relations and global truth valuation operators over state spaces. By replacing rigid element-wise conditions with truth measures $\ddclmd \in [0,1]$, fundamental concepts such as divisibility, primality, and congruence are generalized to capture probabilistic and scenario-based spectral properties of uncertain integers.
	\end{abstract}
	
	\chapter{Preliminaries on Mathematics of Uncertainty}
	
	To ensure this paper is self-contained, we briefly recall key definitions and algebraic structures from the Extended Logic and Mathematics of Uncertainty framework \cite{Tran2026}.
	
	\section{Extended Logic and Truth Valuation}
	Let $(\Omega, \mathcal{F}, \lambda)$ be a measure space representing the scenario/state space of a system. An atomic proposition $P$ possesses a scenario truth function $v(P, \cdot) : \Omega \to \{0, 1\}$.
	
	\begin{definition}[Global Truth Valuation Operator]
		The global truth valuation operator $\ddclmd$ evaluates the truth degree of proposition $P$ over scenario domain $\Omega$:
		\[
		\ddclmd(P) := \frac{\lambda\big(\{\omega \in \Omega : v(P, \omega) = 1\}\big)}{\lambda(\Omega)} \in [0, 1]
		\]
	\end{definition}
	
	\section{Uncertain Numbers and Canonical Representation}
	
	\begin{definition}[Uncertain Number]
		Let $\mathbb{K}$ be an underlying numerical domain (such as a field $\mathbb{R}, \mathbb{C}$ or a commutative ring $\mathbb{Z}$). An uncertain number $A \in \tsbd{K}$ is a quantity whose value belongs to a scenario set $A_s \subseteq \mathbb{K}$.
	\end{definition}
	
	\begin{definition}[Canonical Form]
		An uncertain number $A \in \tsbd{K}$ is represented in canonical form by a pair $(f_A, d_A)$, written $A \sim (f_A, d_A)$, where $d_A \subseteq \mathbb{K}$ is an index domain and $f_A : d_A \to \mathbb{K}$ is a generating function such that $f_A(d_A) = A_s$.
	\end{definition}
	
	\section{Arithmetic Spaces}
	Operations over uncertain numbers interact primarily through two fundamental arithmetic spaces:
	\begin{enumerate}[leftmargin=*]
		\item \textbf{Single-Point Arithmetic Space $(\circ)_1$:} Preserves internal variable correlation and identity across occurrences within an expression:
		\[
		\kgdd{f}(A) = \s{f(x) : x \in A_s}
		\]
		\item \textbf{Minkowski Arithmetic Space $(\circ)_m$:} Models independent scenario interactions and noise propagation over product domains:
		\[
		(A * B)_m = \s{a * b : a \in A_s, b \in B_s}
		\]
		with canonical representation on index space $d_{A * B} = d_A \times d_B$.
	\end{enumerate}
	
	\chapter{Divisibility and Prime Numbers via Weak Relations}
	
	\begin{definition}[Uncertain Integer]
		An element $A \in \tsbd{K}$ is called an uncertain integer if its set of possible scenario values belongs to the set of classical integers $\mathbb{Z}$, i.e., $A_s \subseteq \mathbb{Z}$. The set of all uncertain integers is denoted by $\tsbd{Z}$.
	\end{definition}
	
	\begin{definition}[Weak Divisibility Relation]
		Let $A, B \in \tsbd{Z}$ with canonical forms $A \sim (f_A, d_A)$ and $B \sim (f_B, d_B)$. The weak divisibility proposition $B \mid A$ evaluated on scenario space $\Omega = d_B \times d_A$ has the scenario truth function:
		\[
		v(B \mid A, (t_1, t_2)) = \begin{cases}
			1 & \text{if } f_B(t_1) \mid f_A(t_2) \text{ in } \mathbb{Z}, \\
			0 & \text{otherwise.}
		\end{cases}
		\]
		The global truth valuation of the divisibility relation is given by:
		\[
		\ddclmd(B \mid A) := \frac{\lambda\big(\{(t_1, t_2) \in d_B \times d_A : f_B(t_1) \mid f_A(t_2)\}\big)}{\lambda(d_B \times d_A)}
		\]
		When $\ddclmd(B \mid A) = \gamma$, we write $(B \mid A)_\gamma$. If $\gamma = 1$, $A$ is strictly divisible by $B$.
	\end{definition}
	
	\begin{theorem}[Weak Modus Ponens for Divisibility Transitivity]
		Let $A, B, C \in \tsbd{Z}$. The truth measure of weak divisibility satisfies the inference lower bound:
		\[
		\ddclmd(C \mid A) \ge \ddclmd(C \mid B \land B \mid A)
		\]
		In particular, if $(C \mid B)_1$ and $(B \mid A)_1$, then $(C \mid A)_1$.
	\end{theorem}
	
	\begin{proof}
		Let $A \sim (f_A, d_A)$, $B \sim (f_B, d_B)$, and $C \sim (f_C, d_C)$ be canonical forms. Consider the product state space $\Omega = d_C \times d_B \times d_A$ equipped with product measure $\mu = \lambda_C \times \lambda_B \times \lambda_A$.
		
		Define state subsets of $\Omega$:
		\begin{align*}
			\Omega_{C \mid B} &= \{(t_1, t_2, t_3) \in \Omega : f_C(t_1) \mid f_B(t_2)\}, \\
			\Omega_{B \mid A} &= \{(t_1, t_2, t_3) \in \Omega : f_B(t_2) \mid f_A(t_3)\}, \\
			\Omega_{C \mid A} &= \{(t_1, t_2, t_3) \in \Omega : f_C(t_1) \mid f_A(t_3)\}.
		\end{align*}
		
		By classical transitivity of divisibility in $\mathbb{Z}$, for any state $(t_1, t_2, t_3) \in \Omega$, if $f_C(t_1) \mid f_B(t_2)$ and $f_B(t_2) \mid f_A(t_3)$, then $f_C(t_1) \mid f_A(t_3)$.
		
		Therefore, $\Omega_{C \mid B} \cap \Omega_{B \mid A} \subseteq \Omega_{C \mid A}$. Applying monotonicity of measure $\mu$:
		\[
		\mu(\Omega_{C \mid A}) \ge \mu(\Omega_{C \mid B} \cap \Omega_{B \mid A})
		\]
		Dividing by $\mu(\Omega)$ yields:
		\[
		\ddclmd(C \mid A) \ge \ddclmd(C \mid B \land B \mid A)
		\]
		When $\ddclmd(C \mid B) = 1$ and $\ddclmd(B \mid A) = 1$, we have $\mu(\Omega_{C \mid B}) = \mu(\Omega)$ and $\mu(\Omega_{B \mid A}) = \mu(\Omega)$, so $\mu(\Omega_{C \mid B} \cap \Omega_{B \mid A}) = \mu(\Omega)$, proving $\ddclmd(C \mid A) = 1$.
	\end{proof}
	
	\begin{definition}[Weak Prime Number]
		An uncertain integer $P \in \tsbd{Z}$ with canonical form $P \sim (f_P, d_P)$ is called an uncertain prime number with truth degree $\gamma = \ddclmd(\text{IsPrime}(P))$, where:
		\[
		\ddclmd(\text{IsPrime}(P)) := \frac{\lambda\big(\{t \in d_P : f_P(t) \text{ is a classical prime in } \mathbb{Z}\}\big)}{\lambda(d_P)}
		\]
		If $\gamma = 1$, $P$ is called a strictly prime uncertain number.
	\end{definition}
	
	\begin{definition}[Weak Congruence Relation]
		Let $A, B, N \in \tsbd{Z}$ with canonical forms on $d_A, d_B, d_N$. The weak congruence proposition $A \equiv B \pmod N$ evaluated on state space $\Omega = d_A \times d_B \times d_N$ has truth measure:
		\[
		\ddclmd(A \equiv B \pmod N) := \frac{\lambda\big(\{(t_1, t_2, t_3) \in \Omega : f_A(t_1) \equiv f_B(t_2) \pmod{f_N(t_3)}\}\big)}{\lambda(\Omega)}
		\]
		We write $(A \equiv B \pmod N)_\gamma$ where $\gamma = \ddclmd(A \equiv B \pmod N)$.
	\end{definition}
	
	\begin{proposition}[Properties of Weak Congruence]
		Let $A, B, C, N \in \tsbd{Z}$. Weak congruence satisfies:
		\begin{enumerate}[label=$\mathrm{(\roman*)}$]
			\item \textbf{Symmetry:} $\ddclmd(A \equiv B \pmod N) = \ddclmd(B \equiv A \pmod N)$.
			\item \textbf{Transitivity Bound:} 
			\[
			\ddclmd(A \equiv C \pmod N) \ge \ddclmd(A \equiv B \pmod N) + \ddclmd(B \equiv C \pmod N) - 1
			\]
		\end{enumerate}
	\end{proposition}
	
	\begin{proof}
		$\mathrm{(i)}$ Follows directly from symmetry of classical congruence $a \equiv b \pmod n \iff b \equiv a \pmod n$.
		
		$\mathrm{(ii)}$ Let $\Omega = d_A \times d_B \times d_C \times d_N$. Define measurable sets:
		\begin{align*}
			E_{AB} &= \{(t_1, t_2, t_3, t_4) \in \Omega : f_A(t_1) \equiv f_B(t_2) \pmod{f_N(t_4)}\}, \\
			E_{BC} &= \{(t_1, t_2, t_3, t_4) \in \Omega : f_B(t_2) \equiv f_C(t_3) \pmod{f_N(t_4)}\}, \\
			E_{AC} &= \{(t_1, t_2, t_3, t_4) \in \Omega : f_A(t_1) \equiv f_C(t_3) \pmod{f_N(t_4)}\}.
		\end{align*}
		By classical transitivity of congruence, $E_{AB} \cap E_{BC} \subseteq E_{AC}$. Applying Bonferroni's inequality for measures:
		\[
		\lambda(E_{AC}) \ge \lambda(E_{AB} \cap E_{BC}) \ge \lambda(E_{AB}) + \lambda(E_{BC}) - \lambda(\Omega)
		\]
		Dividing by $\lambda(\Omega)$ establishes the truth measure inequality.
	\end{proof}
	
	\begin{proposition}[Compatibility with Algebraic Operations]
		Let $A, B, C, D, N \in \tsbd{Z}$. In Minkowski space $(\circ)_m$:
		\[
		\ddclmd(A + C \equiv B + D \pmod N) \ge \ddclmd(A \equiv B \pmod N \land C \equiv D \pmod N)
		\]
		\[
		\ddclmd(A \cdot C \equiv B \cdot D \pmod N) \ge \ddclmd(A \equiv B \pmod N \land C \equiv D \pmod N)
		\]
	\end{proposition}
	
	\begin{proof}
		Follows from product measure construction over state domain $d_A \times d_B \times d_C \times d_D \times d_N$ and classical algebraic compatibility of congruences $a \equiv b \pmod n \land c \equiv d \pmod n \implies a+c \equiv b+d \pmod n$ and $a \cdot c \equiv b \cdot d \pmod n$.
	\end{proof}
	
	\begin{example}[Weak Spectrum Congruence Evaluation]
		Let $N = \s{5}$, $A = \s{3, 8}$, and $B = \s{18, 23, 24}$.
		
		The scenario index space is $d_A \times d_B \times d_N$ of size $2 \times 3 \times 1 = 6$ pairs.
		We check classical congruence $a \equiv b \pmod 5$ for each pair $(a, b) \in A \times B$:
		\begin{itemize}
			\item $(3, 18): 3 \equiv 18 \pmod 5$ (True)
			\item $(3, 23): 3 \equiv 23 \pmod 5$ (True)
			\item $(3, 24): 3 \not\equiv 24 \pmod 5$ (False)
			\item $(8, 18): 8 \equiv 18 \pmod 5$ (True)
			\item $(8, 23): 8 \equiv 23 \pmod 5$ (True)
			\item $(8, 24): 8 \not\equiv 24 \pmod 5$ (False)
		\end{itemize}
		The number of satisfying scenario states is 4 out of 6. Therefore, the truth measure of weak congruence is:
		\[
		\ddclmd(A \equiv B \pmod N) = \frac{4}{6} \approx 0.67
		\]
		Thus, we write $\gtcl{A \equiv B \pmod N}{0.67}$.
	\end{example}
	
	\begin{thebibliography}{99}
		\bibitem{Tran2026}
		Tran, M. C. (2026). \textit{Extended logic and mathematics of uncertainty}. Zenodo. DOI:
		\href{https://doi.org/10.5281/zenodo.22031881}{10.5281/zenodo.22031881}
	\end{thebibliography}
\end{document}